\section{Introduction to the LGP-21}
\subsection{Number representation}
\subsubsection*{Introduction to fractional fixed-point architecture}
The LGP-21's fractional fixed-point number representation is very different from more modern computing platforms, though it is not unique; many early computers used it, including the IBM 701\cite{ibm701manual}, the ILLAC I\cite{illiac1954}, the Elliot 405\cite{elliott405}, the RPC-4000\cite{rpc4000manual}, and the LGP-21's predecessor, the LGP-30\cite{lgp30programmingmanual}. This is known as the IAC architecture, and was first proposed by John von Neumann for EDSAC\cite{vonneumann1945first}.

Fixed-point fractional representation treats all numbers as fractions lying strictly between $-1$ and $+1$. As von Neumann pointed out\cite{burks1946preliminary}, the multiplication of 2 $n$-bit integers results in a value that may require $2n$ bits to represent it, which could easily overflow a register, but the multiplication of two \emph{fractions} less than 1 always results in a number no larger than, and usually smaller than, both multiplicands (e.g., $0.5 \times 0.5 = 0.25$). The answer still requires $2n$ bits to represent, but the extra bits end up in the \emph{least} significant digits, preventing overflow from occurring, and making a good compromise between accurate representation and complexity.

This design uses a combined register pair to contain the $2n$ bits; on the LGP-21, the accumulator receives the most significant half, and a hidden multiplication register not available to the programmer gets the least significant half. The LGP-21 simply throws the multiplier register part away. These are the least significant digits, so the (unscaled) error introduced is $2^{-31}$ or less.

This advantage in multiplication is offset by a problem it creates with division: fractional multiplication \emph{shrinks} numbers\footnote{The numerators of numbers representable in fractional fixed-point are strictly smaller than their denominators, so the product of the numerators must always be smaller than the product of the denominators, resulting in a value that is smaller}, but fractional division (potentially) \emph{expands} them.

In a fractional fixed-point architecture, dividing two numbers requires that the absolute magnitude of the divisor has to be \emph{strictly greater than} the absolute magnitude of the dividend.$$\left| \text{Divisor} \right| > \left| \text{Dividend} \right|$$Getting this wrong means that the theoretical quotient is $\ge 1.0$, and such a value overflows because the architecture maintains the \emph{hardware} binary point just after the sign bit\footnote{The LGP-21 also does not have the concept of an exception; the best it can do is set the overflow indicator. This makes division particularly sensitive to scaling factors.}.

Von Neumann wrote that ``\emph{[t]he programmer must see to it} [emphasis mine - ED] that the left-shift [scale factor] is sufficiently large to ensure this inequality.''\cite{burks1946preliminary} The cognitive load of remembering where the \emph{implied} binary point is (via the scaling factor) vs. the \emph{physical} binary point (implemented by the hardware) is placed solely on the programmer.

Modern embedded systems actually still use similar notation (referred to as \emph{Q-notation} or \emph{fixed-point scaling}) to define this manual placement of the implicit binary point within a computer word, and in fact, the original German version of this document refers to the scaling factor as $q$ at various points.

\subsubsection*{LGP-21 Fractional Fixed-point Implementation}
\label{sec:fractional-fixed-point}
The LGP-21 treats all 31-bit words strictly as fractions ranging from $-1$ to $+1$, with the physical binary point fixed immediately to the right of the (leftmost) sign bit. To work with integers or mixed numbers, the programmer must mentally shift this binary point to the right by $n$ bits.

This original version of this manual is inconsistent as to how this is notated (most likely because of multiple authors). As noted, the scaling factor is referred to as $q$, but there are multiple alternate notations for it. Some contributors use $q$ = $n$. Others prefer a shorthand $@n$ notation. Still others simply refer to ``a $q$ of 9'' (e.g.).

For example, a value at $q=3$ allocates 3 bits for the integer portion of a number (allowing values up to $7.999999$ to be represented) and the remaining 27 bits for the fraction. As noted above, the programmer is \emph{entirely responsible} for tracking this scale factor across operations, and must manually shift, multiply, or divide intermediate values as needed to maintain the desired scale factor.

I have retained the melange of different notations for now, but I plan to standardize the usage as $q=n$ for scale factors.

Regardless of how it is notated, the scaling factor represent's the \emph{programmer's} division of the 31 bits of the LGP-21's word into a \emph{logical} integer portion and fractional portion. The hardware binary point remains fixed, but the programmer's binary point can shift. In general, the subroutines that had specific $q$ requirements state them, and note when the $q$ of a result will not match that of original value. 

\subsubsection*{Scale Factors During Arithmetic Operations}
Addition and subtraction require that the $q$ factors of both operands be identical before the operation is executed. The LGP-21 treats every 31-bit word as a fixed-point fraction with the hardware binary point just to the right of the sign bit. If you add two numbers with different $q$ factors without shifting them so their $q$ factors match, the machine will simply perform the operation requested on the two numbers, yielding an invalid result.

    For example, if you attempt to add $1.0$ scaled at $q=1$ to $1.0$ scaled at $q=4$ without normalizing them first:
    \begin{center}
    \begin{tabular}{rrc}
        $1.0 \text{ at } q=1 \rightarrow$ & \texttt{0 100000...} & (Hardware sees $2^{29}$) \\
        $+ \quad 1.0 \text{ at } q=4 \rightarrow$ & \texttt{0 000100...} & (Hardware sees $2^{26}$) \\ \hline
        \text{Result} $\rightarrow$ & \texttt{0 100100...} & (Hardware sees $2^{29} + 2^{26}$)
    \end{tabular}
    \end{center}
Obviously, the answer is not $2.0$, as it should be. The machine will not throw an error, trigger an overflow, or warn you; it will simply calculate perfectly represented garbage because the \emph{logical} binary point in the two numbers is in \emph{two different places}.

Multiplication and division also change the scale factor. Multiplication \emph{adds} $q$ values, division \emph{subtracts} them: a $q=4$ value times a $q=5$ value gives a $q=9$ result. The bit width of the word doesn't change, so the further the result's $q$ drifts from where you want it, the more least significant figures have been lost.

Division has the extra restriction noted earlier: the divisor's magnitude must exceed the dividend's. Equivalently, if subtracting $q_\text{divisor}$ from $q_\text{dividend}$ gives a negative result, the operation overflows and the high-order bits are gone for good. Dividing a $q=1$ value by a $q=9$ value formally yields $q=-8$ --- i.e.\ the top 8, \emph{most} significant bits of the answer were thrown away.
 
It is vital to remember that the LGP-21 hardware itself does not know about scale factors and does not check them or track them, and that this is absolutely, 100\% \emph{your job}.

\subsection{LGP-21 coding conventions}
\subsubsection{Address notation}
Subroutine addresses in the original text are written as \texttt{L} or, for offsets into program code,  \texttt{L+n}. 

Addresses in the main program are written as $\alpha$ or $\alpha+n$.

The LGP-21 is too primitive to have anything like a linkage editor, so it's the programmer/operator's responsibility to use the LGP-21's Program Input Routine\footnote{or Rhys Weatherly's \texttt{lgp21-tools} assembler, which can include and relocate subroutines.} to load subroutine code and automatically relocate it as it's loaded. 

The usage of \texttt{L} for the base address of code is to remind the reader that they will need to relocate these library functions by hand and use the new base addresses
in their code to make subroutine calls work properly. $\alpha\pm$n notation is a reminder that these instructions can live pretty much anywhere that they're not intruding on other loaded code.

\subsubsection{Calling sequences}
Because the LGP-21 libraries were written by multiple programmers, and there was no oversight body to enforce a One True Way of doing things, there is a Cambrian explosion of alternatives. Very capable programmers not only coded their algorithms, but cleverly spaced their storage so that the data was under the disk's read head when an instruction wanted to access memory.\footnote{See the story of Mel\cite{storyofmel}, who originally wrote for the LGP-30; his code was famous for using instructions as data, and exploiting the weird corners of the instruction set to do things like have loops with no exit which nevertheless eventually exit.}

Fortunately, none of those programming tricks are used here\footnote{As far as I know. Heaven knows what's actually lurking in the library code itself.}, but the ``let a million flowers bloom'' philosophy does show itself in the multiple ways there are to call subroutines.

\label{subsec:standard-call1}
\subsubsection*{Accumulator-only parameters}
One calling convention streamlines the call as much as possible by loading the parameter into the accumulator, setting the return address in a \texttt{U} instruction in the subroutine at an agreed-upon offset, and then jumping to the subroutine:
\begin{center}
\begin{tabularx}{\textwidth}{lllX}
\textbf{Location} & \textbf{Instruction} & \textbf{Address} & \textbf{Description} \\
$\alpha-1$      & \texttt{B}            & \texttt{N}       & Bring $N$ into the accumulator. \\
$\alpha$      & \texttt{R}            & \texttt{L+\textit{n}}    & Store the return address in the subroutine. \\
$\alpha+1$      & \texttt{U}            & \texttt{L}       & Unconditional branch to the subroutine entry. \\
\end{tabularx}
\end{center}

In this example a \texttt{B} loads the argument into the accumulator, but any instruction that gets the appropriate value into the accumulator is just as valid. The \texttt{R} saves the return address ($\alpha+2$) in the subroutine, and the main program's \texttt{U} to the start of the subroutine calls it. Execution of the main program resumes just after the \texttt{U} to \texttt{L} at $\alpha+2$.

Here we see the effect of the LGP-21's limited instruction set: there's no such thing as an indirect jump or a stack. Calling a subroutine means that code in that subroutine \emph{must} be modified for a jump back to the caller to work: specifically, a \texttt{U} (unconditional jump) instruction in the subroutine must be altered to set the return address. A second call to the subroutine before it returns will absolutely clobber any previous return address; there's nowhere to put it. Programmers must ensure that they carefully manage the hierarchy of subroutine calls to prevent a second call to a subroutine, directly or indirectly, before it returns.

\subsubsection*{Inlined parameters}
\label{subsec:standard-call2}
A second calling convention, with longer, in-memory parameter lists, looks like this:
\begin{center}
\begin{tabularx}{\textwidth}{lllX}
\textbf{Location} & \textbf{Instruction} & \textbf{Address} & \textbf{Description} \\
$\alpha - 1$      & \texttt{E}            & \texttt{0000}    & Exit from the Interpretive System. \\
$\alpha$          & \texttt{R}            & \texttt{L}       & Modify subroutine instruction with key word address. \\
$\alpha + 1$      & \texttt{U}            & \texttt{L}       & Unconditional branch to subroutine entry. \\
$\alpha + 2$      & \dots                 & \dots            & A fixed number of parameter words \\
$\alpha + n$      & ---                   & ---              & Return to main program. \\
\end{tabularx}
\end{center}
\noindent
This call looks exceedingly strange if you're used to the previous calling convention. If we were doing exactly the same operation as before (altering a \texttt{U} instruction to set the return, then jumping to the subroutine), we'd just jump right back, so that can't be right, and it isn't.

This convention uses the \texttt{R} instruction to load the address of the \emph{parameters} and store it into a \texttt{B} instruction at the start of the subroutine instead. Jumping to the subroutine loads the address of the parameters ($\alpha+2$) into the accumulator, allowing the subroutine to add an offset to this address to set it to $\alpha+n$, which can then be stored locally in a \texttt{U} instruction to do the actual return to the caller, and to have a pointer to the parameters as well!

\subsubsection*{Deep Wizardry: Floating-Point Parameters}

A third, even more astounding calling sequence passes a two-word floating-point\footnote{While the LGP-21 lacks floating-point \emph{hardware}, it relies entirely on a software-driven floating-point math \emph{interpreter}. See Section~\ref{H1-11.0}.} number to a subroutine as an argument.

\begin{center}
\begin{tabularx}{\textwidth}{lllX}
\textbf{Location} & \textbf{Instruction} & \textbf{Address} & \textbf{Description} \\
$\alpha-1$  & \texttt{B} & \texttt{N}        & Bring word 1 of floating-point number $N$. \\
$\alpha$    & \texttt{R} & \texttt{L+offset} & Store the return address ($\alpha+2$) inside the subroutine. \\
$\alpha+1$  & \texttt{U} & \texttt{L}        & Unconditional branch to the subroutine entry. \\
\end{tabularx}
\end{center}

Now wait a minute, I hear you say. A floating-point number is \emph{two} words long. That cannot possibly fit in the single-word hardware accumulator—what kind of trickery is this?

You're right. The full floating-point number doesn't get loaded; only its first word does. But the \texttt{B} instruction itself at cell $\alpha-1$ now contains the exact piece of metadata the subroutine needs: \emph{the starting memory address \texttt{N} of the floating-point block}\footnote{My reaction was ``You are \emph{kidding} me'' with a significant expletive titre.}. 

This explains why the manual repeatedly states that ``loading this value into the accumulator by any means'' is acceptable for some subroutines, but does \emph{not} say that for this one\footnote{Gun, meet foot.}. The \texttt{B} here is essential to the program working. It relies on the code in the subroutine that's modified by the \texttt{R} instruction to \emph{find and read its caller}:
\begin{enumerate}
    \item The subroutine picks up the return address ($\alpha+2$) deposited by the \texttt{R} instruction into cell \texttt{L+offset}.
    \item It subtracts 3 from that value to locate cell $\alpha-1$ (the \texttt{B} instruction).
	\item It modifies the next instruction, a \texttt{B}, to set its address to $\alpha-1$.
    \item The modified \texttt{B} instruction loads the literal 32-bit instruction word stored at $\alpha-1$...that contains the address of the floating-point number.
    \item It strips away the \texttt{B} opcode bits using a bitmask, leaving behind only the raw address \texttt{N}.
	\item Now the subroutine has the address of the first of the two words making up the floating-point number, and can proceed.
\end{enumerate}
The value loaded into the accumulator at step $\alpha-1$ is \emph{entirely ignored}. The \texttt{B} instruction's sole purpose in this calling sequence is to serve as an \emph{executable static data pointer}, which the subroutine finds and dissects to find word 2 (\texttt{N+1}) of the floating-point parameter.

\newpage
